\section{The structure of a research paper}
A scientific paper reports one study in a fixed order: the \emph{abstract} (the whole in 200 words), the \emph{introduction} (the question and why it matters), the \emph{methods} (what was done, in enough detail to repeat it), the \emph{results} (what was found, with the figures and tables), the \emph{discussion} (what it means and where it fails), the \emph{conclusions} and the \emph{references}. The figures carry the results: a good figure can be understood with its caption alone, and the caption says what is plotted, what the panels (a), (b) and (c) are, and what the symbols and colours mean.

\section{How the six papers fit together}
The papers report one project from six sides. The first states the question (what happens to the water stored in the soil when it is released) and answers with the rate-competition law. The second audits the model (does the balance close, what is the term that is not attributed) and gives a compact law of evapotranspiration. The third presents the software that applies the results (the engines and their verification). The fourth asks about the sensor and the trigger (the reach, the fate of a millimetre, the dose). The fifth is about the workbench and the hand-over from a simulation to a decision, and the sixth is about the fields tools and the plants. Read in this order they go from physics to model to tool to use.

\section{Provenance and reproducibility}
A result is \kw{reproducible} if another person, given the data and the code, obtains the same numbers. For this reason every figure of the study is made by a script, every script reads the table of runs, and the table is archived (Zenodo, with a persistent identifier). The \kw{provenance} of a number is the chain from the figure back to the equation, the script, the table and the Daisy run that made it, which the catalogue of figures records. A figure that cannot be traced in this way is an illustration, not a result.

\section{Highlights and graphical abstracts}
A \kw{highlight} is a single sentence of a main finding; a set of three to five gives the reader the content in a minute. A \kw{graphical abstract} is a single picture that shows the question and the answer without text; it is the first thing most readers see, and its aim is to say what the paper is about, not to give the proof. The graphical abstracts of the papers are collected at the end of this chapter.
